% !TeX root = ../papers/nc-p3.tex
\section{滑らかな有理四次点グラフの排除}
\label{sec:rational-quartic-exclusion}
\begin{theorem}\label{thm:rational-quartic-exclusion}
標準有限Frobenius候補で四辺の左右射影ファイバーが空か一点なら、
点グラフが滑らかな整有理四次曲線とそれに交わらない有限スキームの和となることはない。
\end{theorem}
\begin{proof}[二項形式と三割線鉛筆]
$W=H^0(\mathbb P^1,\mathcal O(1))$、$S_n=\operatorname{Sym}^n W$ と書く。
曲線の四次埋込みは非退化である（滑らかな平面四次の種数は3）。
各位置の四切断 $V_i\subset S_4$ の定める中心を $o_i\in\mathbb P(S_4^*)$ とする。
基底 $s,t$ に対して $c_l=o_i(s^{4-l}t^l)$、$C_i=(c_{u+v})_{0\le u,v\le2}$ とする。
このHankel行列は可逆である。特異なら $o_i$ はある二次式 $q$ に対して $qS_2$ を消し、
長さ2の因子の張る割線に乗ってしまう。三次鉛筆
\[
 U_i=\{f\in S_3:o_i(fW)=0\}
\]
は二次元である。$S_3\to W^*$ の階数が2に達しなければ $C_i$ の行も従属する。
共通一次因子 $l$ があれば二つの独立二次式が $C_i$ に関して $lW$ に直交してしまうので
この鉛筆は基点自由である。従って $W\otimes U_i\to V_i$ は同型となる
（互いに素な二つの三次式に一次のsyzygyはない）。
$0\ne f\in U_i$ の零因子は空間三割線に入る。実際、$fW$ の二平面がそこを通る。

隣接する鉛筆は $U_i\cap U_j=0$（以下 $j=i+1,k=i+2$）を満たす。
共通元があれば同じ三割線因子があり、系\ref{cor:quintic-same-center-type}の長さ3の制限と同様、
二割線の積には関係式が高々一本しか残らず、余分な曲線が生じる。
この相補性により $V_i\otimes U_j\to S_7$ は同型である。
鉛筆の基底 $u,v$ との積の核は $(-gv,gu)$、$g\in W$ に限られ、
$gU_j\subset V_i$ は $g=0$ を強制する。実際、$U_j\to W^*$、$f\mapsto o_i(f\,{-})$ は同型である。
さらに $W$ を掛けて $V_i\otimes V_j\to S_8$ は全射で、全核 $R'_{ij}$ は七次元となる。
実際の六関係はその中の超平面 $R_{ij}=\ker\lambda$ である。

\smallskip
\noindent\textit{七関係の二次syzygy束。}
$M_i=\ker(V_i\otimes\mathcal O\to\mathcal O(4))$ とする。
$W\otimes U_i$ の二段の評価列で
$0\to\mathcal O(-1)^2\to M_i\to\mathcal O(-2)\to0$ を得る。
拡大は分裂する。点 $z$ で後ろの文字を評価すると
\[
 0\longrightarrow B\longrightarrow R'_{ij}\otimes\mathcal O
   \longrightarrow M_i(4)\longrightarrow0.
\]
実際、$R_{ij}$ の評価だけで階数3である（$j_j(z)$ の前点が一点）。
$B$ は階数4、次数 $-8$ で $H^0(B(1))=0$ である。
最後の消滅は、$sA+tB'=0$、$A,B'\in R'_{ij}$ なら二つのベクトルが $(tP,-sP)$ で
$P\in V_i\otimes U_j$、かつ $P$ の積が零となることから従う。
よって $B=F\otimes\mathcal O(-2)$、$\dim F=4$。
その包含を二次成分の七行四列行列 $C(z)$ と書く。
行空間 $U\subset F^*\otimes S_2$ は七次元である。定数行の従属は
$H^0(M_i(4)^*)$ に入るが、この空間は零である。
$Q(z)=\lambda C(z)\in F^*\otimes S_2$ はどの $z$ でも零にならない。
もし零なら $B_z\subset R_{ij}$ となり、$R_{ij}$ の評価階数が2に落ちる。

三文字で最初に $R'_{ij}\otimes S_4$ を取り、後ろ二切断の積を零にすると
上の列で $H^0(B(4))=F\otimes S_2$ を得る。最後の切断を $V_k$ に限ると
\[
 K'=(R'_{ij}\otimes V_k)\cap(V_i\otimes R'_{jk})
   =\{v\in F\otimes S_2:\ o_k(Cv)=0\}.
\]
$C_k$ が非退化なので七つの行が課す条件は独立、$\dim K'=12-7=5$ である。
その中に実際の四次元三次交叉があり、$\lambda$ で最初の二位置を縮約した像は高々一次元となる。
すなわち $v\mapsto Qv\in V_k$（二次式の積）の $K'$ での像は高々一本の四次式 $h$ が張る。

\smallskip
\noindent\textit{長さ二の因子上の階数条件。}
$S_2$ 上の点評価列を $e_z=(s(z)^2,s(z)t(z),t(z)^2)^t$ とし、
$q_z=C_k^{-1}e_z\in S_2$ と置く（局所で座標を自明化する）。
$h$ が非零ならその零スキーム上で、像が零なら全直線上で、
\[
                       Q(z)\otimes q_z\in U
\]
が成り立つ。これは $K'$ の零化空間が $(1\otimes C_k)U$ であり、
$Qv$ の点評価が $Q(z)\otimes e_z$ による縮約であることによる。
以下この包含を満たす局所基環上で考える。$Q(z)$ の成分の一つは単元なので、
ある原始的な（剰余体で非零な）行 $\mu$ に対し、独立変数 $x$ について
$\mu C(x)=q_z(x)Q(z)$ となる。

$q_z(x)$ の切る相対次数2の因子を $D$ とする。$D$ 上で $\mu$ は $B$ を消すため、
$M_i(4)|_D$ の双対切断 $\eta$ に降りる。$\eta$ の $R'_{ij}$ の各定数関係での値は定数である。
局所基環上で $D$ の代数は基底 $1,u$ を取り、$u$ の係数を取る汎関数は定数を消し
完全な乗法対合を与える（単位二次式で割った一変数環の表示を取れる）。
これを $\eta$ に合成した原始的な汎関数は $R'_{ij}\to H^0(D,M_i(4)|_D)$ の像を消す。
原始性は、$\mu$ が剰余体で非零でありこの対合が完全なことによる。
この余核は
\[
                  S_6/(V_i\cdot P),\qquad
 P=\ker\bigl(S_2\xrightarrow{\,o_j(q_z\,{-})\,}\mathcal O\bigr)
\]
に等しい。実際、$R'_{ij}\to V_i\otimes V_j\to S_8$ の全射列を $D$ 上の評価列に写す。
中央の評価は全射で、その核は $V_i\otimes q_zP$、右端の核は $q_z S_6$。
蛇の補題で $q_z$ を除けば表示を得る。従って $V_i\otimes P\to S_6$ の七次小行列は全て零になる。

この最後の階数低下の方程式を $P\subset S_2$ の変化する射影平面上で求める。
$P$ が基点自由なら、基底 $u,v$ の積の核は $(-gv,gu)$、$g\in S_2$ で
$o_i(gu)=o_i(gv)=0$ を満たす一次元である。従って積の階数は7。
共通一次因子 $l$ があれば $P=lW$ で、積は $lS_5$ 全体となる。
実際 $W V_i$ は次元 $8-\dim U_i=6$ である。ゆえに台は、$P$ の消滅共変ベクトルが
点評価となる滑らかな円錐曲線である。小行列はこの曲線を被約に切ることも確認できる。
$l(p)=0$、$w(p)\ne0$ とし、基底を $lw,l^2$ から $lw,l^2+\delta w^2$ に動かす。
$g\in U_i$ を $g(p)\ne0$ に選べば、零次のsyzygy $(-gl,gw)$ の一次の像は
$gw^3$ で、余核 $S_6/lS_5$ に非零である。階数6の単元小行列で消去すると、
ある七次小行列の横微分が非零となる。全小行列が円錐曲線で零なので、
各点でこの曲線の被約局所方程式を生成する。

\smallskip
\noindent\textit{Hankel行列の恒等式。}
上の共変ベクトルは $C_jq_z$ であるから、その第一成分と第三成分の積から第二成分の二乗を引いた四次式
\[
 b(z)=(C_jC_k^{-1}e_z)_0(C_jC_k^{-1}e_z)_2-(C_jC_k^{-1}e_z)_1^2
\]
は $h$ の零スキーム上で零である（像が零なら恒等零）。
一方、$A=C_j,B_0=C_k$ とし両者の五つのHankel係数を $a_m,b_m$ と書くと
\[
 (\det B_0)o_k(b)=
 (A\operatorname{adj}(B_0)A)_{02}-(A\operatorname{adj}(B_0)A)_{11}
 =
 \det\!\begin{pmatrix}
 a_0&a_1&a_2&a_3\\a_1&a_2&a_3&a_4\\
 b_0&b_1&b_2&b_3\\b_1&b_2&b_3&b_4
 \end{pmatrix}.
\]
第一の等式は $o_k(e_z e_z^t)=B_0$、第二は展開である。
右辺は $U_j\cap U_k=0$ により非零。従って $b(z)$ は非零で $V_k$ に属さない。
これは像が零の場合に直ちに矛盾し、非零の場合にも同次数の $b$ が $h$ の零因子を含むため
$h\in V_k$ のスカラー倍にならねばならず矛盾する。
\end{proof}

\begin{corollary}
点ファイバーを持つ滑らかな整四次曲線と孤立残余の枝では、
非退化な曲線なら楕円四次しか残らない。平面四次なら種数3なので全位置で平面に入り、
平面の積における三つの一般双一次式の完全交叉の片側次数3に反する。
従って全枝で曲線は楕円四次、孤立部の長さは4である。
\end{corollary}
